\section{Finite moment algebra and the response input}\label{s:moments}

We use finite product sums with explicitly specified moments through degree
four. Pins are arbitrary unit vectors on the original legs; deleting a
site creates no new pin. All degree bounds below count the inserted
observable together with the interaction factors.

\begin{definition}\label{d:moment}
Write \(\mathbb S^2=\{a\in\mathbb R^3:|a|=1\}\), put
\(\phi=(1+\sqrt5)/2\), \(r=\sqrt{\phi+2}\), and set
\[
 D=\frac1r\{(0,\pm1,\pm\phi),(\pm1,\pm\phi,0),(\pm\phi,0,\pm1)\},
 \qquad
 \mathcal M(f)=\frac1{12}\sum_{a\in D}f(a).
\]
All indicated signs are independent. Define finite probability laws
\(\lambda_{\mathrm U}=\frac1{12}\sum_{a\in D}\delta_a\) and
\(\lambda_{\mathrm D}=\frac35\lambda_{\mathrm U}+\frac25\delta_0\).
Products of these functionals always mean finite sums.
For plane sites, \(d\) denotes ambient \(\ell^1\) distance and
\(\operatorname{dist}_\infty\) denotes ambient \(\ell^\infty\) distance.
Distances to an empty set are \(+\infty\); a product over an empty site
set has its one empty configuration. For finite probabilities, use
\[
 d_{\rm TV}(p,p')=\tfrac12\sum_z|p(z)-p'(z)|,\qquad
 \operatorname{Var}_\mu f=\mu(|f-\mu f|^2),\qquad
 d_H(n,n')=\sum_x\mathbf1_{\{n_x\ne n'_x\}} .
\]
\end{definition}

For a vector-valued polynomial, the functional is applied to each
coordinate. In particular the expressions for means below are finite
vector sums.

\begin{lemma}\label{l:moments}
The set \(D\) has twelve distinct unit vectors and is antipodal.
For \(i,j,k,l\in\{1,2,3\}\),
\[
 \begin{gathered}
 \mathcal M(1)=1,\qquad \mathcal M(n_i)=\mathcal M(n_in_jn_k)=0,\qquad
 \mathcal M(n_in_j)=\delta_{ij}/3,\\
 \mathcal M(n_in_jn_kn_l)=
 (\delta_{ij}\delta_{kl}+\delta_{ik}\delta_{jl}+\delta_{il}\delta_{jk})/15.
 \end{gathered}
\]
For every polynomial \(p\) of total degree at most four and every
\(R\in\mathrm O(3)\), \(\mathcal M(p\circ R)=\mathcal M(p)\).
The same covariance holds for \(\lambda_{\mathrm D}\); its constant
moment is one and its moments of positive degree at most four are
\(3/5\) times those displayed above.
These assertions hold site by site for polynomials of separate degree
at most four in a finite product of spins.
\end{lemma}

\begin{proof}
The identities \(\phi^2=\phi+1\), \(r^2=\phi+2=1+\phi^2\) show that
every displayed vector has norm one. Each vector has precisely one zero
coordinate; its location specifies one of the three displayed groups,
and its two nonzero signs specify one of four distinct vectors within
that group. Thus there are twelve distinct vectors. Reversing both signs
gives the antipodal vector.

The set is preserved by changing any individual coordinate sign and by
cyclically permuting the coordinates. A monomial with any odd coordinate
exponent consequently has zero average. For each coordinate \(i\), and
distinct coordinates \(i,j\), direct summation over the three groups gives
\[
 \mathcal M(n_i^2)=\frac{1+\phi^2}{3r^2}=\frac13,\qquad
 \mathcal M(n_i^4)=\frac{1+\phi^4}{3r^4}=\frac15,\qquad
 \mathcal M(n_i^2n_j^2)=\frac{\phi^2}{3r^4}=\frac1{15}.
\]
Here \(r^4=5\phi^2\) and \(1+\phi^4=3\phi^2\).
These cases, the constant, and the vanishing odd-exponent cases exhaust
the monomials through degree four. For four indices, the all-equal case
has three nonzero Kronecker products, the two-equal-pairs case has one,
and every other case has zero; this gives the stated tensor.

For \(R\in\mathrm O(3)\), substituting
\((Rn)_i=\sum_aR_{ia}n_a\) into the second tensor gives
\(\sum_aR_{ia}R_{ja}/3=\delta_{ij}/3\).
Substitution into the fourth tensor gives the three products of the
same contractions, divided by \(15\). The first and third tensors
remain zero, and the constant remains one. Linearity proves covariance
for every polynomial in the stated degree range.

For \(\lambda_{\mathrm D}\), a positive-degree monomial vanishes at
zero, so its moment is \(3/5\) times the corresponding moment of
\(\mathcal M\); the constant moment is \(3/5+2/5=1\).
Zero is fixed by \(R\), so covariance follows as well.
Finally a polynomial of separate degree at most four is a finite sum
of products of single-site monomials of degrees at most four.
Its product functional factors on every such product. Applying the
single-site conclusions to each factor proves the last assertion,
including different orthogonal substitutions at different sites.
\end{proof}

\begin{definition}\label{d:pinmodel}
For finite \(G\subset\mathbb Z^2\), let
\[
 \mathcal L(G)=\{(x,z):x\in G,\ z\notin G,\ d(x,z)=1\},\qquad
 S(G)=\{x\in G:d(x,\mathbb Z^2\setminus G)\le2\}.
\]
A pin assignment is \(\tau\in(\mathbb S^2)^{\mathcal L(G)}\).
For \(Y\subset G\), write \(D_Y=G\setminus Y\) and define
\[
 Z^\tau_{D_Y}[f]=
 \sum_{n\in D^{D_Y}}12^{-|D_Y|}f(n)
 \prod_{\substack{\{x,y\}\subset D_Y\\d(x,y)=1}}(1+n_x\cdot n_y)
 \prod_{\substack{(x,z)\in\mathcal L(G)\\x\in D_Y}}
       (1+\tau_{xz}\cdot n_x).
\]
The free sum \(Z^0_{D_Y}[f]\) omits all leg factors.
Write \(Z^\bullet_{D_Y}=Z^\bullet_{D_Y}[1]\),
\(\mathbb E^\bullet_{D_Y}f=Z^\bullet_{D_Y}[f]/Z^\bullet_{D_Y}\),
and \(\Psi_{D_Y}(\tau)=Z^\tau_{D_Y}/Z^0_{D_Y}\).
Here and below \(\bullet\) is either \(\tau\) or \(0\).
The empty partition sum is one. There are no new factors on legs into \(Y\).
Different original legs may carry unrelated unit pins, even if their
outside endpoints coincide.
\end{definition}

The argument \(f\) is a function of the surviving spins. In the
partition sum every surviving site has at most four factors: each
factor uses one of its four original lattice directions.

\begin{lemma}\label{l:finitefacts}
Consider a finite network with site laws \(\lambda_{\mathrm U}\) or
\(\lambda_{\mathrm D}\), retained distinct-endpoint bonds
\(1+n_x\cdot n_y\), and retained original unit-pin factors
\(1+\tau_\ell\cdot n_x\), with at most four factors incident to each site.
Its pinned and free partition sums are strictly positive, including
the empty network. Non-neighbors are conditionally independent given
the other spins on every conditioning event of positive probability.
If \(f\) is a product of scalar products of spins and the product of
\(f\) with the bond factors has separate degree at most four at every
site, then \(Z^0[f]\ge0\), hence \(\mathbb E^0 f\ge0\).
All normalized spins have norm at most one.
\end{lemma}

\begin{proof}
Every site law has finite support in the closed unit ball, and every
bond or pin factor lies in \([0,2]\). Thus every sum is finite and
nonnegative. To prove strict positivity, order the nonremoved sites
as \(x_1,\ldots,x_m\), and choose at each site a letter in the
twelve-point set \(D\). Such a letter has positive mass under either
law, namely \(1/12\) or \(1/20\).

At stage \(i\), each pin at \(x_i\) excludes at most its antipodal
letter, and each bond to an already chosen site excludes at most the
antipode of that site's letter. There are at most four incident
factors in total, so at most four of the twelve letters are excluded.
Choose any remaining letter. A bond to a later site will be made
strictly positive when that site's letter is chosen. At the end every
retained factor is strictly positive, and the product site mass is
positive. This single summand proves strict positivity. The same
argument with pins omitted proves positivity of the free sum; the
empty network has sum one.

For conditional independence, let \(A\) be any set of sites with no
retained bond between its distinct members, and fix all spins off
\(A\) at a configuration of positive marginal probability. The
unnormalized conditional weight has the form
\[
 C\prod_{x\in A}
 \left(\lambda_{\mathrm{st}(x)}(n_x)
       \prod_{\substack{\text{retained factors}\\\text{incident to }x}}
                     (1+n_x\cdot b)\right).
\]
Here each \(b\) is a fixed neighboring spin or pin, and factors
wholly off \(A\) are included in \(C\). Positive marginal probability
implies \(C>0\) and that every local sum in this product is positive.
The conditional normalizer is \(C\) times the product of those sums.
Dividing therefore gives a product of the single-site conditional
laws. Taking \(A\) to consist of two non-neighbors gives the stated
assertion.

For free positivity expand each retained factor \(1+n_x\cdot n_y\)
into its two terms, and expand every scalar product, including those
in \(f\), as \(\sum_{i=1}^3n_{x,i}n_{y,i}\). Every coefficient in
these finite expansions is nonnegative. For any resulting choice
of coordinate indices, the sum over spins is a product of
single-site moment tensors. The assumed separate degree bound
allows only degrees \(0,1,2,3,4\).
By Lemma~\ref{l:moments}, a degree-one or degree-three tensor is zero,
and the tensors of degrees zero, two and four are sums of products
of Kronecker deltas with nonnegative coefficients.
The same is true for \(\lambda_{\mathrm D}\), with the extra
positive factor \(3/5\) at each site of positive even degree.
Consequently every term after these substitutions is nonnegative.
This also covers scalar products with coincident endpoints, counting
both appearances in the degree. Thus \(Z^0[f]\ge0\); division by the
strictly positive \(Z^0\) proves the normalized assertion.
The supports give \(|n_x|\le1\) pointwise, and hence the same bound
for any conditional or unconditional mean.
\end{proof}

\begin{lemma}\label{l:35}
For \(0\le k\le3\) and \(|a_j|\le1\), the probability law proportional to
\(\prod_{j=1}^k(1+n\cdot a_j)\lambda_{\mathrm U}(\mathrm dn)\)
has mean of norm at most \(3/5\). The same bound holds with
\(\lambda_{\mathrm D}\) in place of \(\lambda_{\mathrm U}\).
Moreover,
\[
 \left\{
 \frac{\mathcal M\!\left(n\prod_{j=1}^3(1+n\cdot a_j)\right)}
      {\mathcal M\!\left(\prod_{j=1}^3(1+n\cdot a_j)\right)}
       :a_1,a_2,a_3\in\mathbb S^2
 \right\}=\{b\in\mathbb R^3:|b|\le3/5\}.
\]
Consequently an unpinned site having at most three present factors
has conditional mean of norm at most \(3/5\), on every positive
conditioning fiber.
\end{lemma}

\begin{proof}
First take three unit vectors \(a,b,c\), and put
\[
 p=a\cdot b+a\cdot c+b\cdot c,\qquad
 A=\mathcal M\bigl((1+n\cdot a)(1+n\cdot b)(1+n\cdot c)\bigr).
\]
The second moment gives \(A=1+p/3\). Since
\(0\le|a+b+c|^2=3+2p\), we have \(A\ge1/2\).
The second and fourth moments give
\[
 15\,\mathcal M\bigl(n(1+n\cdot a)(1+n\cdot b)(1+n\cdot c)\bigr)
 =V,\qquad
 V=5(a+b+c)+a(b\cdot c)+b(a\cdot c)+c(a\cdot b).
\]
Thus the mean is \(V/(5(3+p))\). Set
\(X=1-b\cdot c\), \(Y=1-a\cdot c\), \(Z=1-a\cdot b\);
each lies in \([0,2]\). Expanding the scalar products in \(|V|^2\)
and using \(|a|^2=|b|^2=|c|^2=1\) gives the polynomial identity
\[
 (9+3p)^2-|V|^2
 =4\bigl((X-Y)^2+(X-Z)^2+(Y-Z)^2\bigr)+6XYZ.
\]
The right-hand side is nonnegative, and \(9+3p\ge9/2>0\).
Taking square roots gives \(|V|\le3(3+p)\), which proves the
\(3/5\) bound for three unit fields.

We spell out the reduction to at most three fields in the closed
unit ball. For any such field \(a\), choose a unit vector \(e\)
parallel to \(a\) if \(a\ne0\), and any unit vector if \(a=0\).
With \(p_a=(1+|a|)/2\),
\[
 1+n\cdot a=p_a(1+n\cdot e)+(1-p_a)(1-n\cdot e).
\]
An absent factor can likewise be inserted as
\(1=\tfrac12(1+n\cdot e)+\tfrac12(1-n\cdot e)\).
Expand these identities until there are exactly three unit fields
in every term. The result is a finite convex combination
\(W=\sum_\eta t_\eta W_\eta\), with \(t_\eta\ge0\),
\(\sum_\eta t_\eta=1\), and
\(\mathcal M(W_\eta)=A_\eta\ge1/2\). The normalized mean for \(W\)
is the convex combination of the means for \(W_\eta\) with weights
\(t_\eta A_\eta/\sum_\xi t_\xi A_\xi\).
These weights are nonnegative and sum to one, proving the bound
and positivity of its normalizer in every case \(0\le k\le3\).

For the same product \(W\), let \(A=\mathcal M(W)>0\) and let
\(\mu=\mathcal M(nW)/A\). Since \(W(0)=1\), the mean under the
normalized \(\lambda_{\mathrm D}\) law is
\[
 \frac{(3/5)A}{(3/5)A+2/5}\,\mu.
\]
The prefactor belongs to \((0,1)\), so its norm is at most \(3/5\).

It remains to show that every vector in the ball occurs with three
unit fields under \(\mathcal M\). Take
\[
 a=e_1,\quad b_t=\cos t\,e_1+\sin t\,e_2,\quad
 c_t=\cos t\,e_1-\sin t\,e_2,\qquad 0\le t\le2\pi/3.
\]
The displayed rational formula for the mean has denominator
\(5(3+p)\ge15/2\), so it is continuous along this path.
At \(t=0\), the formula gives mean \(3e_1/5\).
At \(t=2\pi/3\), the three pairwise scalar products are \(-1/2\)
and \(a+b_t+c_t=0\), so \(V=0\).
The continuous norm therefore takes every value in \([0,3/5]\).
The moment formulas also show directly that simultaneously replacing
the three fields by \(Ra,Rb,Rc\) replaces the mean by \(R\mu\),
for every \(R\in\mathrm O(3)\).
For any nonzero desired vector, choose a point on the path with
its norm and an orthogonal matrix sending that mean to the desired
vector; the zero vector already occurs at the endpoint.
The preceding norm bound excludes vectors outside the ball.
This covariance argument concerns only the displayed polynomials,
of degrees three and four.

On a positive conditioning fiber at an unpinned site, its incident
factors give precisely one of the laws just considered, with
neighboring spins of norm at most one. Positivity is as in
Lemma~\ref{l:finitefacts}. Applying the corresponding single-site
bound proves the final assertion.
\end{proof}

\begin{lemma}\label{l:exact}
At each site of a finite set choose any positive weighted unit-vector
law whose moments through degree four equal those of \(\mathcal M\).
A polynomial of separate degree at most four has the same product
functional under these laws as under \(\lambda_{\mathrm U}\) at all
sites. The analogous assertion holds after replacing any chosen site
laws by \(3/5\) times that law plus \(2/5\,\delta_0\).
In particular this applies to partition products of factors
\(1\pm n_x\cdot n_y\) and \(1\pm n_x\cdot a\), \(|a|=1\), whenever
each variable occurs at most four times, and to inserted observables
only when their degrees are included in this same bound.
\end{lemma}

\begin{proof}
Write a polynomial on sites \(x_1,\ldots,x_m\) as a finite linear
combination of monomials
\[
 \prod_{j=1}^m\prod_{i=1}^3 n_{x_j,i}^{\,\gamma_{j,i}},
 \qquad \sum_{i=1}^3\gamma_{j,i}\le4\quad(1\le j\le m).
\]
The product functional of such a monomial is the product of its
single-site moments. Each factor agrees with its
\(\lambda_{\mathrm U}\) value by hypothesis. This proves the first
assertion by linearity, including the empty product.
At a replaced site, the constant moment is unchanged and every
positive-degree moment is multiplied by \(3/5\), since that
monomial vanishes at zero. These are exactly the moments of
\(\lambda_{\mathrm D}\). The same monomial argument proves the
assertion for any specified subset of replaced sites.

A bond factor is affine in each distinct endpoint, and a pin factor
is affine in its spin. Expanding their product therefore has the
claimed separate degree bound when each spin has at most four
appearances, with multiplicity. If an observable is inserted, its
degree is added at each site before this test is applied.
\end{proof}

This replacement identity compares summed-out polynomial expressions.
In particular it supplies the partition-function comparisons for
different finite alphabets with these moments.

\begin{definition}[(P)]\label{h:P}
For \(A\ge2\) and \(0<\theta<1\), property \(\mathsf P(A,\theta)\)
means that, for every finite \(G\subset\mathbb Z^2\), every unit-pin
assignment on \(\mathcal L(G)\), every
\(Y\subset G\setminus S(G)\), and all \(u,v\in G\setminus Y\) with
\(d(u,v)=2\) and a common neighbor in \(Y\),
\[
 \left|\mathbb E^\tau_{G\setminus Y}(n_u\cdot n_v)
             -\mathbb E^0_{G\setminus Y}(n_u\cdot n_v)\right|
 \le A\theta^{\min\{d(u,S(G)),d(v,S(G))\}} .
\]
The distances are ambient plane distances to the original \(S(G)\);
removed sites are not added to the boundary.
\end{definition}


